\chapter{Bleeding Gums}

\section*{Definition and Overview}
Bleeding gums is bleeding from the gingiva, the soft tissue that surrounds and supports the teeth. It is most often noticed while brushing or flossing and, in the great majority of cases, is caused by gingivitis -- the reversible inflammation of the gums that represents the earliest stage of gum disease. Bleeding can also signal more advanced disease (periodontitis), or, less commonly, it may be a manifestation of a systemic condition such as a bleeding disorder, vitamin deficiency, or the effects of blood-thinning medication.

\section*{Epidemiology}
Bleeding on brushing is extremely common, and most adults experience it at some point. Gum disease becomes more frequent with age, and periodontitis is among the most widespread chronic conditions in the world, affecting around half of adults to some degree. It is more common in smokers, in people with poorly controlled diabetes, and in those with inadequate oral hygiene. Spontaneous or heavy gum bleeding, or bleeding accompanied by bruising or bleeding elsewhere, is less common and warrants a search for a systemic cause.

\section*{Aetiology and Pathophysiology}
The most important cause is the accumulation of dental plaque, but several other factors can contribute:

\begin{itemize}
    \item \textbf{Dental plaque:} bacteria gathering at the gum line provoke inflammation, swelling, and fragility of the gums (gingivitis)
    \item \textbf{Periodontitis:} untreated inflammation spreads below the gum line and progressively destroys the supporting bone
    \item \textbf{Aggressive brushing or flossing:} mechanical trauma to the gums
    \item \textbf{Blood-thinning medication:} aspirin, clopidogrel, warfarin, or other anticoagulants increase the tendency to bleed
    \item \textbf{Hormonal changes:} pregnancy gingivitis and bleeding at puberty
    \item \textbf{Vitamin deficiency:} particularly vitamin C (scurvy) and vitamin K
    \item \textbf{Blood disorders:} thrombocytopenia, leukaemia, and clotting-factor deficiencies
    \item \textbf{Ill-fitting dentures, crowns, or fillings:} which irritate the gums
\end{itemize}

Plaque bacteria release toxins that inflame the gingival tissues, causing the small blood vessels to dilate, become fragile, and bleed easily. If plaque is not removed, inflammation deepens, and in periodontitis the enzymes of the inflammatory response destroy the periodontal ligament and alveolar bone that anchor the teeth.

\section*{Clinical Features}
The presentation ranges from minor bleeding noticed at cleaning to spontaneous, painless bleeding:

\begin{itemize}
    \item Blood on the toothbrush or in the rinse basin when brushing or flossing
    \item Red, swollen, or tender gums
    \item Gums that feel sore, spongy, or are receding from the teeth
    \item Persistent bad breath or an unpleasant taste
    \item Loose teeth, shifting teeth, or new gaps between teeth in advanced disease
    \item Gums that bleed on their own, without any brushing
\end{itemize}

\section*{Differential Diagnosis}
\begin{itemize}
    \item Gingivitis, with inflammation but intact supporting structures
    \item Periodontitis, with loss of attachment and bone
    \item Trauma from brushing, flossing, or recent dental treatment
    \item Hormonal gingivitis of pregnancy
    \item Side effects of antiplatelet and anticoagulant medication
    \item Vitamin C or vitamin K deficiency
    \item Blood disorders, including thrombocytopenia, leukaemia, and haemophilia
    \item Bleeding from another oral site, such as a tongue or cheek injury, mistaken for gum bleeding
\end{itemize}

\section*{When to Seek Medical Care}
See a dentist for bleeding that persists despite two weeks of improved brushing and flossing. See a doctor if the gums bleed spontaneously, bleed heavily, or if bleeding is accompanied by unexplained bruising, a pinpoint rash, fever, or unusual tiredness -- features that suggest a systemic disorder.

\textbf{Call 000 (or local emergency services) for:}
\begin{itemize}
    \item Gum bleeding that does not stop with firm pressure after several minutes
    \item Heavy bleeding after dental surgery or an extraction that will not settle
    \item Bleeding gums with confusion, a high fever, or other signs of a serious blood disorder
    \item Sudden widespread bruising or bleeding from other sites
\end{itemize}

\section*{Diagnosis and Investigations}
Assessment distinguishes local gum disease from systemic causes of bleeding.

\begin{itemize}
    \item \textbf{Dental examination:} inspection of the gums, gentle probing to measure pocket depths, and assessment of tooth mobility and bone loss
    \item \textbf{Oral hygiene review:} evaluation of brushing and flossing technique
    \item \textbf{Radiographs:} dental X-rays to assess the supporting bone in suspected periodontitis
    \item \textbf{Blood tests:} full blood count, platelet count, clotting studies, and vitamin or iron levels for unexplained or severe bleeding
    \item \textbf{Medication review:} including over-the-counter blood thinners
    \item \textbf{Haematology referral:} when a bleeding disorder is suspected
\end{itemize}

\section*{Management}
Treatment is directed at the cause, with control of inflammation as the central aim:

\begin{itemize}
    \item \textbf{Improved oral hygiene:} regular twice-daily brushing with a soft brush and daily interdental cleaning
    \item \textbf{Professional cleaning:} scaling and polishing to remove plaque and tartar, performed by a dentist or hygienist
    \item \textbf{Periodontal therapy:} root surface debridement (deep cleaning) and, in advanced disease, surgical treatment
    \item \textbf{Follow-up maintenance:} scheduled hygiene visits to prevent recurrence
    \item \textbf{Treatment of the underlying cause:} dose adjustment of anticoagulants on medical advice, vitamin replacement, or specialist management of a blood disorder
    \item \textbf{Softer brushing:} a gentle technique and a soft or extra-soft brush where trauma is contributing
\end{itemize}

\section*{Complications}
\begin{itemize}
    \item Progression of gingivitis to periodontitis with irreversible bone loss
    \item Loosening and loss of teeth
    \item Persistent bad breath and altered taste
    \item Gum infection and periodontal abscess formation
    \item Wider health associations: periodontitis is linked with poorer diabetes control and increased cardiovascular risk
    \item Significant blood loss in the rare cases caused by bleeding disorders
\end{itemize}

\section*{Prognosis and Outcomes}
Bleeding gums caused by gingivitis usually settle within one to two weeks of consistent good oral hygiene and professional cleaning. Periodontitis cannot be reversed, but it can be effectively controlled with ongoing maintenance care so that most teeth are retained for life. Where bleeding reflects a systemic cause, it resolves with treatment of that cause, and the prognosis is that of the underlying condition.

\section*{Prevention and Self-Care}
\begin{itemize}
    \item Brush twice daily with a fluoride toothpaste and a soft-bristled brush
    \item Clean between the teeth daily with floss or interdental brushes
    \item Do not stop cleaning bleeding gums -- gentle cleaning is what makes them healthy, and bleeding should reduce within a week or two
    \item Attend regular dental check-ups and professional cleans
    \item Avoid smoking, which markedly worsens gum disease
    \item Manage diabetes and other chronic conditions carefully
    \item See a dentist promptly for persistent, spontaneous, or heavy gum bleeding
\end{itemize}

\section*{Sources and Further Reading}
The following sources provide reliable, up-to-date information for patients and students of medicine.

\begin{itemize}
    \item NHS. \textit{Overview of gum disease and how to look after your gums.} https://www.nhs.uk/conditions/gum-disease/
    \item Mayo Clinic. \textit{Gingivitis -- symptoms, causes, and treatment.} https://www.mayoclinic.org/diseases-conditions/gingivitis/
    \item MSD Manual. \textit{Gum disease and bleeding gums -- professional edition.} https://www.msdmanuals.com/
    \item MedlinePlus. \textit{Bleeding gums -- patient information.} https://medlineplus.gov/
    \item Centers for Disease Control and Prevention. \textit{Periodontal disease -- information and resources.} https://www.cdc.gov/oral-health/
\end{itemize}
